\section{Conclusion}

We began by asking whether gradient-based interventions could provide single-run robustification against spurious correlations without group labels. Our investigation reveals that before this question can be answered, a prerequisite must be met: gradient subspace methods require sufficient accumulation density to produce meaningful directional measurements.

\subsection{Summary}

Our experiments on Waterbirds with ResNet-50 yielded a surprising null result---spurious and core alignments were both approximately 0.05, statistically indistinguishable. Investigation traced this to measurement apparatus failure: accumulating one gradient per epoch produced a 10-dimensional subspace in a 25-million-parameter space.

Our contributions are methodological:

\begin{enumerate}
    \item We establish minimum requirements for gradient subspace analysis: multi-batch accumulation within early epochs.
    \item We document a reproducible failure case that future gradient-based methods should avoid.
    \item We provide guidance for valid testing of gradient orthogonalization hypotheses.
\end{enumerate}

\subsection{Future Directions}

\textbf{Immediate:} Re-implement gradient accumulation using streaming SVD across all batches during epochs 1-10.

\textbf{Medium-term:} If valid measurement confirms spurious-core separation, implement full PGO mechanism and evaluate against baselines.

\textbf{Long-term:} Extend to other architectures and domains where gradient subspace properties may differ.

\subsection{Closing}

The promise of gradient-based robustification remains unrealized. What we have established is what valid measurement requires. Future work should validate accumulation density before claiming directional findings. The question of whether early gradients capture spurious directions remains open, but now we know how to test it properly.
